\section{Experimental research}\label{sec:experiments}

This chapter describes the experiments conducted to evaluate the proposed system.
The reporting follows the benchmarking rules of Hoefler and Belli~\cite{hoefler2015scientific}.

\subsection{Experimental setup}

The experiments were performed on the following hardware:
\begin{itemize}
  \item CPU: 8-core processor, \qty{3.2}{\giga\hertz};
  \item RAM: 32~GB DDR4;
  \item storage: 512~GB NVMe SSD;
  \item OS: Ubuntu 24.04 LTS.
\end{itemize}

\subsection{Test datasets}

Three datasets of varying sizes were used for evaluation (\cref{tab:datasets}).

% EXAMPLE: table with booktabs and siunitx; thousands are grouped automatically.
\begin{table}[htbp]
  \centering
  \caption{Test datasets}\label{tab:datasets}
  \begin{tabular}{l S[table-format=7.0] S[table-format=4.0] l}
    \toprule
    Dataset & {Records} & {Size (MB)} & Purpose             \\
    \midrule
    Small   & 1000      & 5           & Unit testing        \\
    Medium  & 100000    & 450         & Integration testing \\
    Large   & 1000000   & 4200        & Stress testing      \\
    \bottomrule
  \end{tabular}
\end{table}

\subsection{Performance results}

The system was benchmarked against two baseline implementations.
\Cref{tab:performance} shows the results.
The raw numbers are stored in \path{figures/data/benchmark.csv}.

% EXAMPLE: grouped column headers with \multicolumn and \cmidrule.
% (multirow is avoided: it is not yet compatible with PDF tagging.)
\begin{table}[htbp]
  \centering
  \caption{Performance comparison (processing time in seconds)}\label{tab:performance}
  \begin{tabular}{l S[table-format=1.1] S[table-format=2.1] S[table-format=3.1]}
    \toprule
                      & \multicolumn{3}{c}{Dataset size}           \\
    \cmidrule(lr){2-4}
    Method            & {Small} & {Medium} & {Large} \\
    \midrule
    Baseline A        & 0.8     & 45.2     & 512.0   \\
    Baseline B        & 0.5     & 38.7     & 478.3   \\
    Proposed approach & 0.3     & 22.1     & 287.5   \\
    \bottomrule
  \end{tabular}
\end{table}

The proposed approach reduces processing time by 44--63\% compared to Baseline~A and by 40--43\% compared to Baseline~B.

% EXAMPLE: pgfplots bar chart that reads its data from a CSV file.
\Cref{fig:bar-chart} visualizes the performance comparison across dataset sizes.

\begin{figure}[htbp]
  \centering
  \begin{tikzpicture}[alt={Grouped bar chart of processing time in seconds.
      Small dataset: Baseline A 0.8, Baseline B 0.5, Proposed 0.3. Medium:
      45.2, 38.7, 22.1. Large: 512.0, 478.3, 287.5. The proposed approach is
      fastest for every dataset size.}]
    \begin{axis}[
        ybar,
        bar width=14pt,
        enlarge x limits=0.25,
        legend style={at={(0.5,-0.15)}, anchor=north, legend columns=-1,
          /tikz/every even column/.append style={column sep=1em}},
        ylabel={Time (s)},
        symbolic x coords={Small, Medium, Large},
        xtick=data,
        ymin=0, ymax=600,
        width=12cm, height=8cm,
      ]
      \addplot[fill=red!60]   table[col sep=comma, x=dataset, y=baselineA] {figures/data/benchmark.csv};
      \addplot[fill=blue!60]  table[col sep=comma, x=dataset, y=baselineB] {figures/data/benchmark.csv};
      \addplot[fill=green!60] table[col sep=comma, x=dataset, y=proposed]  {figures/data/benchmark.csv};
      \legend{Baseline A, Baseline B, Proposed}
    \end{axis}
  \end{tikzpicture}
  \caption{Performance comparison across dataset sizes}\label{fig:bar-chart}
\end{figure}

\subsection{Accuracy results}

Accuracy was measured using standard classification metrics~\cite{sokolova2009systematic}:
\begin{equation}\label{eq:accuracy}
  \text{Accuracy} = \frac{TP + TN}{TP + TN + FP + FN} \times 100\%
\end{equation}
where $TP$, $TN$, $FP$, $FN$ denote true positives, true negatives, false positives, and false negatives respectively.

% EXAMPLE: aligned equations with extra vertical space.
The precision, recall, and $F_1$-score are defined as:
\begin{align}
  \text{Precision} & = \frac{TP}{TP + FP} \label{eq:precision} \\[6pt]
  \text{Recall}    & = \frac{TP}{TP + FN} \label{eq:recall}    \\[6pt]
  F_1              & = 2 \cdot \frac{\text{Precision} \cdot \text{Recall}}{\text{Precision} + \text{Recall}} \label{eq:f1}
\end{align}

The proposed approach achieved an accuracy of 94.2\%, compared to 89.5\% for Baseline~A and 91.1\% for Baseline~B (\cref{tab:metrics}).

\begin{table}[htbp]
  \centering
  \caption{Detailed classification metrics (\%)}\label{tab:metrics}
  \begin{tabular}{l *{4}{S[table-format=2.1]}}
    \toprule
    Method            & {Accuracy} & {Precision} & {Recall} & {$F_1$-score} \\
    \midrule
    Baseline A        & 89.5       & 87.3        & 88.1     & 87.7          \\
    Baseline B        & 91.1       & 90.2        & 89.8     & 90.0          \\
    Proposed approach & 94.2       & 93.5        & 93.8     & 93.6          \\
    \bottomrule
  \end{tabular}
\end{table}

\subsection{Scalability analysis}

The system scales linearly with dataset size, as evidenced by the following relationship:
\begin{equation}\label{eq:scalability}
  T(n) = a \cdot n + b
\end{equation}
where $T(n)$ is the processing time, $n$ is the number of records, and $a$, $b$ are constants determined through regression analysis ($a = \num{2.21e-4}$, $b = \num{0.12}$).

% EXAMPLE: pgfplots line plot from a CSV file.
\Cref{fig:scalability-plot} shows the scalability curve.

\begin{figure}[htbp]
  \centering
  \begin{tikzpicture}[alt={Line chart of processing time in seconds against
      records in hundreds of thousands. The proposed approach grows linearly
      from 0.3 seconds to 176.8 seconds at 800 thousand records; Baseline B
      grows faster, reaching 260 seconds.}]
    \begin{axis}[
        xlabel={Records ($\times 10^5$)},
        ylabel={Time (s)},
        xmin=0, xmax=8.5,
        ymin=0, ymax=300,
        legend pos=north west,
        grid=major,
        width=12cm, height=8cm,
      ]
      \addplot[thick, green!50!black, mark=*] table[col sep=comma, x=records, y=proposed] {figures/data/scalability.csv};
      \addlegendentry{Proposed}
      \addplot[thick, blue!70, dashed, mark=square*] table[col sep=comma, x=records, y=baselineB] {figures/data/scalability.csv};
      \addlegendentry{Baseline B}
    \end{axis}
  \end{tikzpicture}
  \caption{Scalability: processing time versus dataset size}\label{fig:scalability-plot}
\end{figure}
